\section{Window continuity and the parity of a simple minimum}
\label{sec:spectral-window-continuity}

\begin{theorem}[AC-S46: norm-resolvent continuity and open simple-window sets]
\label{AC-S46}
Identify \(L^2(0,\ell)\) with \(L^2(0,1)\) by
\[
                   (U_\ell f)(x)=\sqrt\ell\,f(\ell x).
                                                               \tag{WC.1}
\]
Then \(U_\ell A_\ell U_\ell^{-1}\) depends continuously
on \(\ell>0\) in norm-resolvent sense.  In particular the
two lowest eigenvalues are continuous functions of \(\ell\),
and the multiple-minimum set \(\mathcal M\) in (GS.1)
is relatively closed.  Together with AC-S45, it is a closed
countable subset of \((0,\infty)\).

The simple-minimum set is open and dense.  On each of its
connected components, the parity of its real lowest line is
constant.  Both the simple-even and simple-odd sets are open;
any change of parity must pass through \(\mathcal M\).

There exists \(\delta>0\) such that the actual minimum is
simple and even throughout
\[
       0<\ell<\log r_0+\delta,
       \qquad\hbox{equivalently}\qquad
       1<\lambda<\sqrt{r_0}\,e^{\delta/2},
       \quad r_0^3-r_0-1=0.                              \tag{WC.2}
\]
Thus the proved simple-even interval extends beyond the
ordinary-cone threshold in AC-S28.  This argument gives no
explicit positive value of \(\delta\), no nonnegativity
assertion for its eigenfunction beyond that threshold, and
no proof that every simple minimum is even.
\end{theorem}
\begin{proof}
\proofstep{1. The logarithmic part on a fixed interval.}
Fix \(\ell_0>0\) and restrict \(\ell\) to a compact
neighborhood of \(\ell_0\) in \((0,\infty)\).
After (WC.1), the full-line logarithmic multiplier is
\[
                  g_\ell(\xi)=\tfrac12
                            \log(1+\xi^2/\ell^2).
\]
Direct differentiation with respect to \(\log\ell\) gives
\[
 \left|\frac{\partial}{\partial\log\ell}g_\ell(\xi)\right|
                =\frac{\xi^2}{\ell^2+\xi^2}\leq1,
 \qquad
 \|g_\ell-g_{\ell_0}\|_\infty
                        \leq|\log(\ell/\ell_0)|.        \tag{WC.3}
\]
Its zero-extension part on \((0,1)\) therefore differs
from the fixed part \(G_0\) at \(\ell_0\) by a bounded
self-adjoint operator \(C_\ell\), with
\(\|C_\ell\|\to0\).
This is an equality of closed forms and of their associated
operators, so their operator domains agree.

\proofstep{2. The bounded part and entering prime atoms.}
The integrable kernel \(k\) in (RN.4) becomes
\(\ell k(\ell(x-y))\).
The maps \(\ell\mapsto\ell k(\ell\,\cdot)\) are
continuous in \(L^1(\mathbb R)\): prove this first for
compact continuous kernels, and then approximate \(k\)
in \(L^1\), using that dilation preserves its \(L^1\) norm.
Thus these compressed convolution operators vary in operator
norm.  The polar kernel becomes
\(2\ell\cosh(\ell(x-y)/2)\) on \((0,1)^2\), and varies
there continuously in Hilbert--Schmidt norm.
The scalar term is constant.

Only finitely many prime powers can enter in the chosen
neighborhood of \(\ell_0\).  Include all of them throughout
that neighborhood.  Each gives fixed scalar multiples of
\[
                   P S_{\pm\log m/\ell}E,                \tag{WC.4}
\]
where \(E\) is zero extension from \((0,1)\) and \(P\)
is restriction back to it.
Whole-line translations are strongly continuous on \(L^2\),
so (WC.4) and their adjoints vary strongly continuously.
When \(|\log m/\ell|\geq1\), the compression is zero
as an \(L^2\) operator.  This includes equality, so the
same formula handles the precise lengths at which a prime
atom enters; no norm continuity of sharp translations is
being asserted.

Consequently the scaled operator has the form
\[
                  \widetilde A_\ell
                    =G_0+C_\ell+B_\ell,                 \tag{WC.5}
\]
where \(B_\ell\) is uniformly bounded and strongly
continuous, and each operator is self-adjoint on the fixed
domain \(\operatorname{Dom}(G_0)\).
These operators have a common lower bound on the chosen
parameter neighborhood.

\proofstep{3. Compactness converts strong variation into resolvent norm variation.}
Choose \(c\) so large that \(\widetilde A_\ell+c\geq1\)
throughout the neighborhood, and set
\(R_\ell=(\widetilde A_\ell+c)^{-1}\).
The resolvent identity is
\[
 R_\ell-R_{\ell_0}
  =-R_\ell(C_\ell+B_\ell-B_{\ell_0})R_{\ell_0}.          \tag{WC.6}
\]
Here \(R_{\ell_0}\) is compact and \(\|R_\ell\|\leq1\).
A uniformly bounded strongly convergent sequence of operators
converges uniformly on a compact set: approximate that set
by a finite net and use the uniform bound on the remaining
errors.  Equivalently, approximate \(R_{\ell_0}\) in norm
by finite-rank operators.
It follows that the right side of (WC.6) tends to zero in
operator norm.  This proves norm-resolvent continuity;
the same identity at any fixed nonreal spectral parameter
gives the usual equivalent definition directly.

The compact positive operators \(R_\ell\) now vary in norm.
The min--max characterization bounds the change of each
ordered eigenvalue by \(\|R_\ell-R_{\ell_0}\|\).
Their first two eigenvalues are positive, and inversion
\(\mu\mapsto\mu^{-1}-c\) therefore proves continuity of
the two lowest eigenvalues of \(A_\ell\).
In particular their equality set \(\mathcal M\) is closed
relative to \((0,\infty)\), and its complement is open.
AC-S45 makes that complement dense and of full measure.

\proofstep{4. Parity and the extension beyond the cone threshold.}
On the fixed space \(L^2(0,1)\), the reflection
\(Jf(x)=f(1-x)\) is independent of \(\ell\) and commutes
with every scaled operator.  Near a simple minimum the
rank-one spectral projector \(P_\ell\) varies in norm.
For example, integrate the resolvent of the norm-continuous
compact operator \(R_\ell\) around a fixed small circle
isolating its largest simple eigenvalue; the Neumann series
gives uniform convergence on that circle.
The scalar \(\operatorname{Tr}(JP_\ell)\) is continuous
and takes only the values \(1\) and \(-1\).
The trace continuity follows also from
\(\operatorname{rank}(P_\ell-P_{\ell_0})\leq2\)
and norm convergence.
Thus parity is locally constant on the simple set and
constant on each connected component.

At \(\ell=\log r_0\), AC-S28 proves a simple even
minimum, including the equality case of its kernel bound.
Openness and local constancy give a neighborhood of that
point with the same property.  Combining its right-hand
part with the full interval already covered by AC-S28
gives (WC.2) for some \(\delta>0\).
Nothing in this continuity argument preserves the
nonnegative cone after its kernel threshold, so no
pointwise positivity conclusion is added there.
\end{proof}
